\section{Method}
\label{sec:method}

\subsection{Problem Setup}

Let $x \in \mathcal{X}$ denote a street-view image and
$a \in \mathcal{A}$ a perceptual attribute (e.g.\ safety). We take the
explanatory atom to be a \emph{lever intervention}
\begin{equation}
  l = (c, s, d, \tau),
\end{equation}
where $c$ is a lever concept (e.g.\ greenery, graffiti), $s$ is the
grounded scene support, $d$ the intervention direction (add, remove,
repair), and $\tau$ a constrained edit template.

For image $x$, let $L(x)$ denote the scene-specific candidate set
instantiated from an editor-feasible intervention vocabulary. Our method identifies the
subset that can be realised as valid edits:
\begin{equation}
  V(x) = \{ l_i \in L(x) \mid \exists\, \tilde{x}_i \text{ passing the validity audit} \}.
\end{equation}
Each lever is tested independently against the unedited original; edits
are never composed sequentially.

\subsection{Phase 1: Lever Candidate Construction}

We define an \emph{editor-feasible intervention vocabulary}: a bounded
set of lever concepts chosen to satisfy three criteria:
\textbf{(i)}~\emph{empirical grounding} in the urban perception
literature~\cite{jacobs1961,newman1972,li2015greenery,jiang2018brokenwindows,portnov2020lighting},
\textbf{(ii)}~\emph{semantic distinctness} so concepts remain
interpretable, and \textbf{(iii)}~\emph{prompt-only editability}
without requiring segmentation masks or major structural changes.
The full concept list is given in Appendix~\ref{sec:intervention_vocabulary}.

For each image, a VLM planner grounds applicable vocabulary concepts to
visible scene regions, producing scene-specific candidates such as
\emph{remove graffiti from this wall} or \emph{add modest greenery near
this entrance}. These grounded candidates form $L(x)$.

\subsection{Phase 2: Bounded Stochastic Intervention}

For each candidate $l_i$, a prompt-only image generator $G$ produces
\begin{equation}
  x'_{i,j} = G(x;\, c_i, s_i, d_i, \tau_i, \varphi_j),
\end{equation}
where $\varphi_j$ indexes stochastic draws under a bounded retry budget
$T$, with $G$ being a prompt-only image editing model.

Each candidate is subjected to an automated validity audit performed by
a vision--language model,
evaluating: \textbf{(1)}~same-place preservation,
\textbf{(2)}~locality to the intended support, \textbf{(3)}~realism,
and \textbf{(4)}~plausibility of the intended lever. If at least one
draw passes within $T$ attempts, the first accepted edit $\tilde{x}_i$
is retained and $l_i$ enters $V(x)$. Only edits passing all four
criteria enter the retained set, so auxiliary shifts are measured only
on edits with confirmed semantic and spatial integrity.
We report valid-edit coverage
$\mathrm{Coverage}(x) = |V(x)| / |L(x)|$ as a property of edit
realisability under bounded prompt control rather than an urban
perception signal.

\subsection{Phase 3: Auxiliary Scoring}

For each accepted edit $\tilde{x}_i$, we compute
\begin{equation}
  \Delta_i = f_a(\tilde{x}_i) - f_a(x),
\end{equation}
where $f_a$ is a ViT-B/16 perception model pretrained on the MIT Place
Pulse~2.0 pairwise safety dataset~\cite{hou2024globalstreetscapes},
which outputs a continuous score on a 0--10 scale (higher = safer).
This is strictly an auxiliary proxy, not the main estimand; $f_a$ has
not been validated on edited images, so $\Delta_i$ should be treated as
a ranking signal rather than a calibrated effect size. The proxy-shortlisted set is
\begin{equation}
  E_{\mathrm{aux}}(x,a) = \{ l_i \in V(x) \mid \Delta_i >
  \theta_{\mathrm{aux}} \},
\end{equation}
where $\theta_{\mathrm{aux}}$ is an exploratory threshold.
Human-grounded pairwise evaluation remains future
work.%
\footnote{Our method is constrained interventional search, not
pixel-level decomposition or causal identification. ``Multiple levers''
means distinct interventions that each pass validity checks
independently; it does not imply additive effects. Locality is audited,
not pixel-guaranteed.}
